\documentclass[12pt,a4paper]{report}

% ===================== PREAMBLE =====================

\usepackage{listings}
\usepackage{xcolor}
\usepackage{adjustbox}

\definecolor{lightgray}{gray}{0.95}

\lstset{
    backgroundcolor=\color{lightgray},
    basicstyle=\ttfamily\small,
    frame=single,
    breaklines=true,
    captionpos=b,
    keepspaces=true,
    showstringspaces=false
}

\usepackage{enumitem}
\setlist[itemize]{noitemsep, topsep=0pt}
\setlist[enumerate]{noitemsep, topsep=0pt}


% ----- Margins -----
\usepackage[a4paper, top=1in, bottom=1in, left=1.25in, right=1.25in]{geometry}

% ----- Times New Roman -----
\usepackage{mathptmx}

% Line spacing
\usepackage{setspace}
\setstretch{1.5}

% Additional packages
\usepackage{graphicx}
\usepackage{titlesec}
\usepackage{hyperref}
\usepackage{float}
\usepackage{caption}
\usepackage{parskip}
\usepackage{multirow}
\usepackage{booktabs}
\usepackage{tikz}
\usetikzlibrary{shapes.geometric, arrows, positioning}

% ----- Page Number at Bottom Center -----
\usepackage{fancyhdr}
\pagestyle{fancy}
\fancyhf{}
\fancyfoot[C]{\thepage}
\renewcommand{\headrulewidth}{0pt}

% ----- Heading Formatting -----
% Main headings - 14pt, bold, uppercase, centered
\titleformat{\section}
{\normalfont\fontsize{14pt}{16pt}\bfseries\centering\MakeUppercase}
{\thesection}{1em}{}

% Subheading – 12pt, bold, sentence case, left
\titleformat{\subsection}
{\normalfont\fontsize{12pt}{14pt}\bfseries}
{\thesubsection}{1em}{}

% Sub-subheading – 12pt, bold italic, left
\titleformat{\subsubsection}
{\normalfont\fontsize{12pt}{14pt}\bfseries\itshape}
{\thesubsubsection}{1em}{}

% ----- Caption Formatting -----
\captionsetup{
    font={bf,11pt},
    labelsep=period
}
\setcounter{secnumdepth}{3}
\renewcommand{\thesection}{\arabic{section}}
\renewcommand{\thesubsection}{\arabic{section}.\arabic{subsection}}


% ----------------- START DOCUMENT -----------------
\begin{document}

% ----------------- TITLE PAGE -----------------
\begin{titlepage}
\thispagestyle{empty}

\begin{center}

% -------- College Logo --------
\includegraphics[width=3cm]{cbit_logo.png}\\[0.4cm]

{\Large \textbf{CHAITANYA BHARATHI INSTITUTE OF TECHNOLOGY}}\\
{\normalsize (Autonomous)}\\[0.4cm]


{\Large \textbf{BIG DATA ANALYTICS}}\\
{\large Assignment -- 1}\\[0.6cm]

{\large \textbf{REPORT ON}}\\[0.3cm]

{\LARGE \textbf{Post-Mortem Analysis of the MapReduce Framework}}\\
{\Large \textbf{Using Large-Scale Taxi Trip Data Analytics}}\\[0.8cm]



\begin{tabular}{c}
\textbf{Submitted by} \\[0.3cm]
Mohammed Faizan Ul Islam \quad (160123737195)\\
Vedith Vanam \quad (160123737207)\\
Shashidhar Nagunuri \quad (160123737317)\end{tabular}\\[0.9cm]

\textbf{Under the Guidance of}\\
Dr. N. Sudhakar Yadav\\[0.6cm]

Department of Information Technology\\
Chaitanya Bharathi Institute of Technology\\[0.6cm]

\textbf{Academic Year: 2025--2026}

\end{center}

\end{titlepage}


% ----------------- CERTIFICATE PAGE -----------------
\newpage

\begin{center}
{\LARGE \textbf{ABSTRACT}}
\end{center}

\addcontentsline{toc}{section}{Abstract}

\vspace{0.5cm}

The rapid growth of urban transportation systems and digital taxi services has resulted
in the continuous generation of large volumes of trip-level data. Each taxi trip record
captures spatial, temporal, operational, and financial information, making such datasets
valuable for large-scale analytical studies. However, traditional single-node processing
techniques are inefficient for handling datasets containing millions of records.

This report presents a comprehensive post-mortem analysis of the Hadoop MapReduce
programming model using large-scale taxi trip data analytics as a case study. The dataset
was processed using Hadoop MapReduce across multiple storage formats including CSV,
Parquet, and Avro. Hadoop Streaming with Python-based mapper and reducer programs
was used for CSV processing, while Java MapReduce implementations were used for
Parquet and Avro datasets.

Multiple MapReduce jobs were executed to analyze pickup locations, drop-off locations,
passenger count distribution, payment type frequency, trip type classification, and trip
frequency distribution across hours of the day. The study evaluates MapReduce
performance with respect to scalability, fault tolerance, execution time, and storage
efficiency under different file formats.

The results demonstrate that storage format significantly influences distributed processing
performance. CSV requires higher parsing overhead, while Parquet provides columnar
compression advantages and Avro offers efficient binary serialization. These findings
provide practical insights into the role of storage formats in improving efficiency and
scalability in distributed big data processing frameworks.

\clearpage


% ----------------- CERTIFICATE PAGE -----------------

\newpage

% ----------------- TOC / LOF / LOT -----------------
\addcontentsline{toc}{section}{TABLE OF CONTENTS}
\tableofcontents
\clearpage
\listoffigures
\addcontentsline{toc}{section}{LIST OF FIGURES}
\clearpage
\listoftables
\addcontentsline{toc}{section}{LIST OF TABLES}
\clearpage

\pagenumbering{arabic}
\setcounter{page}{1}
\setlength{\parindent}{20pt} % Paragraph indentation
\setlength{\parskip}{1.0em} % Space between paragraphs

% ----------------- INTRODUCTION -----------------
\newpage
\pagenumbering{arabic}
\section{Introduction}

The rapid growth of urban transportation systems and digital taxi services has led to the
continuous generation of large volumes of trip-level data. Taxi trip datasets record
detailed information such as pickup and drop-off locations, timestamps, passenger
counts, trip distance, and fare-related attributes. When collected over long periods,
these datasets grow to millions of records and become valuable for understanding
transportation demand and travel behavior.

Analyzing taxi trip data enables transportation authorities and service providers to
identify demand hotspots, study passenger patterns, and evaluate operational
efficiency. However, the large size and complexity of such datasets pose significant
challenges for traditional single-machine data processing techniques. Centralized
systems suffer from limited scalability, high execution latency, and poor fault tolerance
when handling large-scale data.

To address these challenges, distributed data processing frameworks have become
essential. The MapReduce programming model provides a scalable approach for
processing large datasets by dividing computations into map and reduce phases that
can be executed in parallel. Apache Hadoop offers an open-source implementation of
MapReduce along with a distributed storage layer, enabling reliable batch-oriented
analytics on commodity hardware.

In this study, the MapReduce framework is applied to a large taxi trip dataset to evaluate
its effectiveness for batch analytics. The dataset is processed using multiple file formats
including CSV, Parquet, and Avro to compare execution performance and storage
efficiency. Sequential processing is also considered as a baseline for understanding the
benefits of distributed computation.

Multiple MapReduce jobs are executed to analyze different attributes of the dataset,
including pickup location frequency, drop-off location frequency, passenger count
distribution, payment type frequency, trip type classification, and hourly trip distribution.
These analyses enable evaluation of MapReduce performance and execution behavior
under realistic data characteristics.


\subsection{Motivation}

Although MapReduce has been widely adopted for batch-oriented big data processing,
its real-world performance characteristics are best understood through empirical
evaluation. Transportation datasets often exhibit repetitive categorical values, skewed
distributions, and varying aggregation requirements, all of which influence MapReduce
execution.

Different file formats can also influence performance characteristics. CSV requires text
parsing during processing, while Parquet provides columnar storage benefits and Avro
supports efficient binary serialization. Evaluating these formats helps in understanding
how storage structure affects distributed processing efficiency.

The motivation of this study is to analyze the behavior of the MapReduce framework
when applied to a realistic large-scale taxi trip dataset stored in multiple file formats.


\subsection{Problem Statement}

The primary problem addressed in this report is the scalable analysis of large-scale taxi
trip data using the MapReduce programming model. The study seeks to evaluate how
MapReduce performs when processing millions of trip records and executing multiple
aggregation-based analytical workloads with different key distributions.

Additionally, the study compares performance differences when the dataset is stored in
CSV, Parquet, and Avro formats, and contrasts distributed processing with sequential
processing approaches.


\subsection{Objectives of the Study}

The objectives of this study are as follows:

\begin{itemize}
    \item To process large-scale taxi trip data using the MapReduce framework.
    \item To implement multiple MapReduce jobs for analyzing transportation data.
    \item To evaluate MapReduce performance through post-execution analysis.
    \item To compare CSV, Parquet, and Avro file formats in distributed processing.
\end{itemize}


\subsection{Scope of the Report}

The scope of this report is limited to batch-oriented analytics performed using the
Apache Hadoop MapReduce framework. The analysis focuses on descriptive and
aggregation-based processing of taxi trip data using CSV, Parquet, and Avro file formats.

Real-time analytics, predictive modeling, and in-memory processing frameworks are
beyond the scope of this study.


\clearpage
% ----------------- EXISTING SYSTEM -----------------
\section{Background and Related Work}

The rapid growth of data-intensive applications across domains such as transportation,
finance, healthcare, and e-commerce has significantly influenced the evolution of data
processing systems. Early data management solutions were primarily designed to handle
structured data using centralized architectures. Relational database management systems
and traditional batch processing tools were effective for moderate workloads but relied
heavily on vertical scaling to improve performance.

As data volume and complexity increased, the limitations of centralized processing
became increasingly evident. Single-node systems struggled to process large datasets
within acceptable time limits and were highly susceptible to hardware and software
failures. Furthermore, many real-world datasets, including transportation data, exhibit
semi-structured characteristics that are difficult to manage using conventional database
approaches.

These challenges led to the development of distributed data processing models that
emphasize horizontal scalability, parallel execution, and fault tolerance. Distributed
systems allow datasets to be partitioned across multiple machines, enabling simultaneous
processing and improved reliability.

The MapReduce programming model, originally proposed by Google, simplified distributed
processing by dividing computation into two operations—map and reduce. Apache Hadoop
provided an open-source implementation of MapReduce along with the Hadoop
Distributed File System (HDFS), enabling scalable batch-oriented analytics.

In modern big data ecosystems, datasets are commonly stored in formats such as CSV,
Parquet, and Avro. CSV is simple and widely supported but requires text parsing during
processing. Parquet provides columnar storage that improves compression and query
efficiency, while Avro supports binary serialization with schema information. These file
formats influence storage efficiency and processing performance in distributed systems.

Although MapReduce remains effective for large-scale batch processing, newer frameworks
such as Apache Spark have been developed to address limitations related to latency and
iterative processing. Nevertheless, MapReduce continues to be widely used for reliable
large-scale aggregation tasks.
\clearpage
% ----------------------------------
\section{MapReduce Framework Overview}

The MapReduce framework is a distributed programming model designed to process
large-scale datasets across clusters of commodity hardware. It simplifies parallel data
processing by abstracting complexities such as task scheduling, data distribution, and
fault tolerance. This allows developers to build scalable data processing applications
without managing low-level distributed system details.

MapReduce operates in close integration with the Hadoop Distributed File System
(HDFS), which provides reliable and scalable storage for large datasets. Input data is
divided into blocks and distributed across multiple nodes. These blocks may contain
datasets stored in formats such as CSV, Parquet, or Avro, and are processed in parallel
by mapper tasks. Intermediate results are aggregated by reducer tasks to produce the
final output.


\subsection{Hadoop Distributed File System (HDFS)}

HDFS is a distributed storage system designed to store large files across multiple
machines while ensuring reliability and scalability. It follows a master–slave architecture
consisting of a NameNode and multiple DataNodes. The NameNode maintains metadata
such as file names and block locations, while DataNodes store the actual data blocks.

To provide fault tolerance, HDFS replicates each data block across multiple nodes. In the
event of node failure, data can be retrieved from another replica without interrupting
processing. HDFS supports multiple file formats including CSV, Parquet, and Avro.


\subsection{Map Phase}

The map phase is the first stage of execution. Each mapper processes a portion of input
data and transforms it into intermediate key–value pairs. The mapper typically performs
operations such as parsing records, filtering data, and extracting required attributes.

In transportation data analysis, mapper functions extract attributes such as pickup
location, drop-off location, passenger count, payment type, trip type, and timestamp.
Each extracted attribute is emitted as a key associated with a value representing its
occurrence.


\subsection{Shuffle and Sort Phase}

The shuffle and sort phase connects the map and reduce phases. During this stage,
intermediate key–value pairs are redistributed across the cluster so that all values
associated with the same key are grouped together and sent to the appropriate reducer.

This phase involves network communication and disk operations, making it one of the
most resource-intensive stages of MapReduce execution.


\subsection{Reduce Phase}

The reduce phase aggregates grouped intermediate data produced during the shuffle
phase. Each reducer processes a unique key and list of associated values to compute
final results such as counts or summaries.

The aggregated output is written back to HDFS for further analysis.


\subsection{Fault Tolerance and Data Locality}

Fault tolerance is a key feature of MapReduce. If a mapper or reducer task fails, the
framework automatically re-executes the task on another node, ensuring reliable
completion of jobs on datasets stored in CSV, Parquet, or Avro formats.

MapReduce also uses data locality by scheduling tasks on nodes where data is stored,
reducing network transfer and improving execution efficiency.
\clearpage

% ----------------------------------
\section{Case Study Description}

This study focuses on the large-scale analysis of taxi trip data using the MapReduce
framework. Modern urban transportation systems generate extensive trip-level data as
a result of continuous daily operations. Each taxi trip produces a structured record
containing information related to trip origin, destination, passenger details, payment
characteristics, and timestamp information. Due to their high volume and repetitive
nature, taxi trip datasets are well-suited for batch-oriented distributed processing
frameworks such as MapReduce.

The selected case study evaluates the effectiveness of MapReduce for processing and
analyzing large transportation datasets stored in multiple file formats including CSV,
Parquet, and Avro. These workloads are designed to vary in key distribution and
aggregation complexity, enabling detailed observation of MapReduce execution behavior
and performance characteristics.


\subsection{Analytical Problem Definition}

The primary objective of this case study is to extract meaningful insights from a large
taxi trip dataset in a scalable and fault-tolerant manner. Specifically, the study aims to
address the following analytical questions:

\begin{itemize}
    \item Which pickup locations generate the highest number of taxi trips?
    \item Which drop-off locations are most frequently used by passengers?
    \item How are taxi trips distributed based on passenger count?
    \item What payment methods are most commonly used for taxi services?
    \item How are taxi trips classified based on trip type?
    \item How are taxi trips distributed across different hours of the day?
\end{itemize}

Answering these questions using traditional sequential processing approaches becomes
computationally expensive as dataset size increases. Therefore, the MapReduce
framework is used to perform distributed aggregation-based analytics.


\subsection{Dataset Description}

The dataset used in this study consists of over one million taxi trip records collected
during the year 2020. Each record corresponds to an individual taxi trip and includes
attributes related to pickup location, drop-off location, passenger count, payment type,
trip type, and timestamp information.

The dataset is stored in CSV, Parquet, and Avro formats to evaluate how storage
structure influences distributed processing performance. The size and diversity of the
dataset make it suitable for evaluating scalability and efficiency of the MapReduce
framework.


\subsection{Data Record Format}

The dataset is represented in multiple structured formats.

CSV (Comma-Separated Values) stores records in text format where attributes are
separated by commas. It is simple and widely supported but requires parsing during
processing.

Parquet stores data in columnar format, enabling efficient compression and improved
performance for analytical queries that access specific attributes.

Avro stores data in binary row format along with schema information, providing compact
storage and efficient serialization in distributed environments.

Using multiple file formats enables comparison of execution time, storage efficiency,
and processing complexity within Hadoop MapReduce.


\subsection{Rationale for Case Study Selection}

Taxi trip data analytics was selected as the case study due to its practical relevance,
large data volume, and suitability for batch-oriented distributed processing. Transportation
datasets contain repetitive categorical values such as locations, payment types, and trip
categories, making them suitable for aggregation-based analysis using MapReduce.

The dataset characteristics allow evaluation of MapReduce performance under realistic
conditions including skewed key distributions and large data volume. Comparing CSV,
Parquet, and Avro formats also helps in understanding how storage structure affects
processing efficiency in distributed systems.
\clearpage 

% ----------------------------------
\section{Design and Implementation of MapReduce Jobs}

This section describes the design and implementation of the MapReduce jobs developed
for analyzing large-scale taxi trip data. Each analytical task is implemented as an
independent MapReduce job. CSV datasets are processed using Hadoop Streaming with
Python-based mapper and reducer programs, while Parquet and Avro datasets are
processed using Java MapReduce with ParquetInputFormat and AvroKeyInputFormat.

All MapReduce jobs follow a common design pattern. The mapper extracts a relevant
attribute from each input record and emits it as a key with a value of one. A shared
reducer program is used across all jobs to aggregate values associated with each key
and compute final counts. This ensures consistency of logic across CSV, Parquet, and
Avro datasets.


\subsection{Pickup Location Frequency Analysis}

The pickup location frequency analysis identifies areas with the highest taxi demand by
counting the number of trips originating from each pickup location. The mapper extracts
the pickup location identifier from each taxi trip record and emits it as a key with value
one. The reducer aggregates these values to compute the total number of trips for each
pickup location.

This analysis helps in understanding spatial demand patterns and identifying high-traffic
zones. Due to repetitive location identifiers, this workload also allows observation of key
skew and reducer load imbalance during MapReduce execution.


\subsection{Drop-off Location Frequency Analysis}

The drop-off location frequency analysis identifies commonly used destinations in taxi
trips. The mapper extracts the drop-off location identifier from each record and emits it
as a key-value pair with value one. The reducer aggregates these values to compute the
total number of trips ending at each location.

This analysis provides insight into passenger destination trends and allows evaluation of
MapReduce performance under uneven key distribution.


\subsection{Passenger Count Distribution Analysis}

The passenger count distribution analysis examines how taxi trips are distributed based
on the number of passengers per trip. The mapper extracts the passenger count
attribute and emits it as the key. The reducer aggregates values to compute frequency
distribution of passenger counts.

Since the number of distinct passenger count values is limited, this workload results in
relatively balanced reducer execution compared to location-based analyses.


\subsection{Payment Type Frequency Analysis}

The payment type frequency analysis evaluates customer payment preferences by
counting the number of trips associated with each payment method. The mapper extracts
the payment type field and emits it as a key with value one. The reducer aggregates
counts to produce distribution of payment methods.

This workload demonstrates MapReduce performance for datasets with small key space,
resulting in efficient reducer execution.


\subsection{Trip Type Distribution Analysis}

The trip type distribution analysis classifies taxi trips based on service category. The
mapper extracts the trip type attribute and emits it as the key. The reducer aggregates
values to determine distribution of trip types.

This workload provides insight into operational characteristics of taxi services and
demonstrates aggregation on categorical attributes.


\subsection{Hourly Trip Frequency Analysis}

The hourly trip frequency analysis computes the total number of taxi trips occurring in
each hour of the day. The mapper extracts the pickup timestamp from each record,
converts it into hour format (0–23), and emits the hour as the key with value one.

The reducer aggregates values corresponding to each hour to produce the total number
of trips occurring during each hour of the day. This job is executed on CSV, Parquet, and
Avro datasets to compare execution performance across file formats.


\subsection{Common Reducer Implementation}

A single reducer implementation is used across all MapReduce jobs to perform
aggregation. The reducer receives intermediate key–value pairs sorted by key and sums
values associated with each key to compute final counts.

Using a common reducer ensures consistency of aggregation logic across all analytical
workloads and file formats. The reducer outputs results to HDFS, where they are later
retrieved for analysis and comparison.
\clearpage
% ----------------------------------
\section{File Format Based Implementation}

To evaluate the impact of storage format on distributed processing performance,
the same taxi trip dataset was stored and processed using three different formats:
CSV, Parquet, and Avro. The objective was to observe differences in execution
time, storage efficiency, parsing complexity, and MapReduce performance.

\subsection{CSV Format Processing}

The dataset was initially stored in CSV format and processed using Hadoop Streaming
with Python-based mapper and reducer programs. CSV is a row-based text format,
requiring parsing of each record before extracting attributes such as pickup timestamp.

Characteristics observed:
\begin{itemize}
\item Human readable format
\item Requires text parsing for each record
\item Larger file size compared to binary formats
\item Higher I/O overhead
\end{itemize}

\subsection{Parquet Format Processing}

The dataset was converted into Parquet format and processed using Java MapReduce
with ParquetInputFormat. Parquet is a columnar storage format optimized for
analytical queries.

Characteristics observed:
\begin{itemize}
\item Columnar storage improves compression efficiency
\item Faster read performance for analytical workloads
\item Reduced disk I/O
\item Schema-based structure
\end{itemize}

\subsection{Avro Format Processing}

The dataset was stored in Avro format and processed using Java MapReduce with
AvroKeyInputFormat. Avro is a row-based binary serialization format designed for
efficient distributed processing.

Characteristics observed:
\begin{itemize}
\item Compact binary representation
\item Faster serialization and deserialization
\item Embedded schema improves data consistency
\item Suitable for row-based processing
\end{itemize}

\subsection{Objective of Format Comparison}

The same MapReduce logic was applied across all formats to compute trip frequency
per hour. Performance differences were evaluated based on execution time,
storage efficiency, and processing complexity.

\clearpage

\section{Results, Observations, and Post-Mortem Evaluation}

This section presents the results obtained from executing multiple MapReduce jobs on
the large-scale taxi trip dataset. The analytical workloads were executed on datasets
stored in CSV, Parquet, and Avro formats. The outputs generated by each workload are
examined to identify meaningful patterns in the data and evaluate execution behavior of
the MapReduce framework.

The same MapReduce logic was applied across all file formats to ensure consistent
outputs. Differences in performance were observed due to variations in storage structure,
compression, and parsing complexity.


\subsection{Result Summary}

Each MapReduce job produced structured output files stored in the Hadoop Distributed
File System (HDFS). The outputs were retrieved and examined to verify correctness and
analyze aggregate trends in the taxi trip data. Successful execution across CSV,
Parquet, and Avro datasets demonstrates the ability of the MapReduce framework to
process large-scale batch workloads efficiently.

\begin{table}[h]
\centering
\caption{Summary of Taxi Trip Analytical Workloads}
\begin{tabular}{|l|l|l|}
\hline
\textbf{Analysis Type} & \textbf{Mapper Used} & \textbf{Output Directory} \\
\hline
Pickup Location Analysis & mapper\_pu.py & /user/taxi/output\_pu \\
Drop-off Location Analysis & mapper\_do.py & /user/taxi/output\_do \\
Passenger Count Analysis & mapper\_passengers.py & /user/taxi/output\_passenger \\
Payment Type Analysis & mapper\_payment.py & /user/taxi/output\_payment \\
Trip Type Analysis & mapper\_triptype.py & /user/taxi/output\_triptype \\
Hourly Trip Analysis & mapper\_hour.py & /user/taxi/output\_hour \\
\hline
\end{tabular}
\end{table}


\subsection{Execution on CSV, Parquet, and Avro}

MapReduce jobs were executed on datasets stored in CSV, Parquet, and Avro formats.
Execution logs show successful completion of mapper and reducer tasks across all file
formats.

CSV datasets required additional parsing during execution due to text-based storage.
Parquet datasets demonstrated improved disk read efficiency due to columnar storage,
while Avro datasets provided efficient binary serialization with consistent schema
representation.

\begin{figure}[h]
\centering
\includegraphics[width=0.9\textwidth]{csv_execution.png}
\caption{CSV MapReduce Execution}
\end{figure}

\begin{figure}[h]
\centering
\includegraphics[width=0.9\textwidth]{parquet_execution.png}
\caption{Parquet MapReduce Execution}
\end{figure}

\begin{figure}[h]
\centering
\includegraphics[width=0.9\textwidth]{avro_execution.png}
\caption{Avro MapReduce Execution}
\end{figure}


\subsection{Pickup Location Frequency Analysis}

The pickup location frequency analysis computes the total number of taxi trips
originating from each pickup location. The results show that a limited number of
locations account for a large proportion of taxi demand.

This skewed distribution results in uneven reducer workloads during MapReduce
execution across all file formats.

\begin{figure}[h]
\centering
\includegraphics[width=0.9\textwidth]{Picture3.png}
\caption{Pickup Location Frequency Analysis Output}
\end{figure}


\subsection{Drop-off Location Frequency Analysis}

The drop-off location analysis identifies commonly used destinations in taxi trips.
Certain drop-off locations dominate the distribution, indicating concentrated travel
demand patterns.

This workload demonstrates shuffle overhead caused by skewed key distribution in
MapReduce processing.

\begin{figure}[h]
\centering
\includegraphics[width=0.9\textwidth]{Picture5.png}
\caption{Drop-off Location Frequency Analysis Output}
\end{figure}


\subsection{Passenger Count Distribution Analysis}

The passenger count distribution analysis examines how taxi trips are distributed based
on number of passengers per trip. The results indicate that trips with one or two
passengers occur most frequently.

Due to limited distinct values, this workload results in balanced reducer execution.

\begin{figure}[h]
\centering
\includegraphics[width=0.9\textwidth]{Picture7.png}
\caption{Passenger Count Distribution Output}
\end{figure}


\subsection{Payment Type Frequency Analysis}

The payment type frequency analysis evaluates distribution of payment methods used
by passengers. The results show dominant usage of certain payment types.

Balanced key distribution results in efficient reducer execution.

\begin{figure}[h]
\centering
\includegraphics[width=0.9\textwidth]{Picture10.png}
\caption{Payment Type Frequency Distribution Output}
\end{figure}


\subsection{Trip Type Distribution Analysis}

The trip type distribution analysis categorizes taxi trips based on service classification.
Certain trip types account for majority of services.

Limited key variation results in efficient MapReduce aggregation.

\begin{figure}[h]
\centering
\includegraphics[width=0.9\textwidth]{Picture9.png}
\caption{Trip Type Distribution Output}
\end{figure}


\subsection{Hourly Trip Frequency Analysis}

The hourly analysis computes total number of taxi trips occurring in each hour of the
day. The output consists of 24 key-value pairs representing trip frequency distribution
from hour 0 to hour 23.

This workload was executed on CSV, Parquet, and Avro datasets to compare execution
performance across storage formats.

\begin{figure}[h]
\centering
\includegraphics[width=\textwidth]{hour_output.png}
\caption{Hourly Trip Frequency Output}
\end{figure}


\subsection{Comparative Performance Analysis}

Comparison of execution behavior shows that storage format influences processing
efficiency. CSV format requires text parsing which increases processing overhead.
Parquet format improves performance through columnar compression and reduced disk
I/O. Avro format provides efficient binary serialization and schema consistency.

Location-based analyses exhibit higher shuffle overhead due to skewed key distribution,
while categorical analyses such as payment type and trip type demonstrate balanced
reducer workloads.

\begin{table}[H]
\centering
\caption{Execution Time Comparison Across Different Systems}

\begin{tabular}{|c|c|c|c|c|}
\hline

\textbf{System} & \textbf{Sequential} & \textbf{CSV (Hadoop)} & \textbf{Parquet (Hadoop)} & \textbf{Avro (Hadoop)} \\

\hline

System 1 & 9m 23s & 20m 53s & 10m 10s & 5m 30s \\

System 2 & 8m 30s & 18m 30s & 9m 50s & 4m 58s \\

System 3 & 18m 30s & 40m 00s & 22m 14s & 11m 18s \\

\hline

\end{tabular}

\end{table}



\subsection{Comparison with Sequential Processing}

Sequential processing executes operations on a single machine without parallelism.
Processing large datasets sequentially results in higher execution time and limited
scalability.

MapReduce improves performance by dividing data into multiple splits processed in
parallel across distributed nodes. Fault tolerance mechanisms ensure reliable job
completion even when node failures occur.

Distributed processing significantly reduces execution time compared to sequential
processing when handling large datasets.

\begin{table}[H]
\centering
\caption{Comparison of Sequential Processing and Hadoop-based File Formats}

\begin{adjustbox}{width=\textwidth}

\begin{tabular}{|l|c|c|c|c|}
\hline

\textbf{Aspect} & \textbf{Sequential} & \textbf{CSV} & \textbf{Parquet} & \textbf{Avro} \\

\hline

Execution Model & Single node & Distributed & Distributed & Distributed \\

Storage Location & Local FS & HDFS & HDFS & HDFS \\

Parallelism & No & Yes & Yes & Yes \\

Scalability & Low & High & High & High \\

Fault Tolerance & No & Yes & Yes & Yes \\

Large Dataset Handling & Poor & Good & Very Good & Very Good \\

Execution Time & Highest & High & Medium & Lowest \\

Processing Speed & Slow & Moderate & Fast & Fast \\

Storage Efficiency & Low & Medium & High & High \\

Data Format & Text & Row text & Columnar & Binary row \\

Complexity & Low & Medium & Medium & Medium \\

Best Use Case & Small data & General big data & Analytics & Data exchange \\

\hline

\end{tabular}

\end{adjustbox}

\end{table}


\subsection{Post-Mortem Evaluation}

Based on observed execution behavior, MapReduce performs reliably for large-scale
batch aggregation tasks. The framework provides scalability, parallel processing, and
fault tolerance.

However, shuffle and sort phase introduces disk and network overhead, especially for
skewed key distributions. CSV processing requires additional parsing effort, while Parquet
and Avro improve storage efficiency through optimized data representation.

These observations highlight how both processing model and storage format influence
performance in distributed big data environments.
\newpage
\section{Conclusion and Future Scope}

\subsection{Conclusion}

This report presented a comprehensive post-mortem analysis of the Hadoop MapReduce
programming model using large-scale taxi trip data as a case study. Multiple MapReduce
jobs were designed and executed to analyze attributes such as pickup location frequency,
drop-off location frequency, passenger count distribution, payment type frequency, trip
type classification, and hourly trip distribution.

The dataset was processed using CSV, Parquet, and Avro file formats to evaluate how
storage structure influences distributed processing performance. The experimental results
demonstrate that MapReduce efficiently processes large-scale datasets through parallel
execution, fault tolerance, and scalable architecture.

CSV format provides simplicity and compatibility but requires additional parsing due to
its text-based structure. Parquet format improves efficiency through columnar storage
and compression, making it suitable for analytical workloads. Avro format provides
compact binary serialization along with schema support, enabling efficient distributed
data exchange.

Comparison with sequential processing shows that distributed MapReduce significantly
reduces execution time and improves scalability when handling large datasets. The
post-mortem evaluation highlights that storage format and data characteristics influence
shuffle overhead, disk usage, and execution efficiency.


\subsection{Future Scope}

Future work can extend this study by implementing the same analytical workloads using
modern distributed processing frameworks such as Apache Spark to compare execution
performance with MapReduce.

Further analysis can also be performed using larger datasets and multi-node cluster
configurations to evaluate scalability under real-world big data environments. Additional
analytical tasks such as trip duration analysis, fare distribution analysis, and peak demand
prediction can provide deeper insights into transportation patterns.

Visualization tools and dashboards can also be integrated to improve interpretability of
analytical results and support data-driven decision making.
\clearpage
% ----------------------------------
\end{document}